\documentclass[conference]{IEEEtran}
\usepackage[spanish]{babel}
\usepackage[utf8]{inputenc}
\usepackage[T1]{fontenc}
\usepackage{amsmath}
\usepackage{booktabs}
\usepackage{url}

\title{Limites para Ordenar Complejidades Algoritmicas}
\author{\IEEEauthorblockN{Ronquillo Nunez Braulio}
\IEEEauthorblockA{Analisis y Diseno de Algoritmos}}

\begin{document}
\maketitle

\begin{abstract}
El ordenamiento de funciones de complejidad es una tarea basica en analisis de
algoritmos. Para comparar dos costos se puede estudiar el limite del cociente
entre sus funciones. Este reporte desarrolla el criterio por limites, presenta
una jerarquia usual de crecimiento y describe un programa que ordena
expresiones de complejidad de menor a mayor clase asintotica.
\end{abstract}

\begin{IEEEkeywords}
limites, complejidad, Big-O, jerarquia asintotica, algoritmos.
\end{IEEEkeywords}

\section{Introduccion}
Cuando se comparan algoritmos, no basta con observar un caso pequeno. Un
algoritmo con constantes altas puede superar a otro cuando $n$ crece, y un
algoritmo rapido en entradas pequenas puede volverse impracticable si su
crecimiento es exponencial. Por eso se estudian funciones como $\log n$, $n$,
$n\log n$, $n^2$, $2^n$ y $n!$.

NIST define Big-O como una medida teorica del tiempo o memoria necesarios para
un algoritmo segun el tamano de entrada \cite{nist_big_o}. Esta idea se puede
complementar con limites para decidir cuando una funcion domina a otra.

\section{Criterio por Limites}
Sean $f(n)$ y $g(n)$ funciones positivas. Se analiza
\begin{equation}
L=\lim_{n\to\infty}\frac{f(n)}{g(n)}.
\end{equation}
Si $L=0$, entonces $f$ crece estrictamente menos que $g$. Si $L$ es una
constante positiva y finita, ambas funciones son del mismo orden. Si
$L=\infty$, entonces $f$ crece mas que $g$.

La regla de L'Hospital puede ayudar en ciertos cocientes indeterminados, pero
debe aplicarse con cuidado. MathWorld indica que la regla compara limites de
derivadas bajo condiciones especificas y tambien muestra casos donde la
aplicacion mecanica puede fallar por oscilacion \cite{mathworld_lhopital}.

\section{Jerarquia Usual}
Una escala comun de crecimiento es:
\begin{equation}
1 \prec \log n \prec n \prec n\log n \prec n^2 \prec n^3 \prec 2^n \prec n!.
\end{equation}
Esta jerarquia no cubre todas las funciones posibles, pero es suficiente para
clasificar muchas complejidades introductorias. La notacion $\Theta$ formaliza
la idea de misma clase asintotica mediante cotas superior e inferior por
constantes \cite{nist_theta}.

\begin{table}[htbp]
\caption{Clases usadas por el programa}
\centering
\begin{tabular}{ll}
\toprule
Clase & Ejemplo \\
\midrule
Constante & $1$ \\
Logaritmica & $\log n$ \\
Lineal & $n$ \\
Lineal-logaritmica & $n\log n$ \\
Polinomial & $n^2$, $n^3$ \\
Exponencial & $2^n$ \\
Factorial & $n!$ \\
\bottomrule
\end{tabular}
\end{table}

\section{Implementacion}
El programa recibe una lista de expresiones escritas con nombres predefinidos.
Cada expresion se asocia con un rango numerico que representa su posicion en
la jerarquia. Luego se ordena la lista por ese rango. Cuando dos expresiones
comparten rango, se imprimen como equivalentes en la misma clase.

Esta solucion es deliberadamente didactica: en lugar de hacer algebra
simbolica completa, codifica una tabla de clases para practicar el concepto de
dominancia asintotica.

\section{Complejidad}
Si se reciben $k$ expresiones, el costo dominante es el ordenamiento, por lo
que el tiempo es $O(k\log k)$. La memoria adicional es $O(k)$ para almacenar
las expresiones y sus claves de comparacion.

\section{Conclusion}
El analisis por limites da una base formal para ordenar complejidades. Aunque
el programa usa una jerarquia cerrada, reproduce la idea central: comparar el
crecimiento dominante y agrupar funciones equivalentes bajo $\Theta$.

\begin{thebibliography}{00}
\bibitem{nist_big_o} P. E. Black, ``big-O notation,'' Dictionary of Algorithms and Data Structures, NIST, 2019. [Online]. Available: \url{https://www.nist.gov/dads/HTML/bigOnotation.html}
\bibitem{nist_theta} P. E. Black, ``Theta,'' Dictionary of Algorithms and Data Structures, NIST, 2016. [Online]. Available: \url{https://www.nist.gov/dads/HTML/theta.html}
\bibitem{mathworld_lhopital} E. W. Weisstein, ``L'Hospital's Rule,'' MathWorld--A Wolfram Resource. [Online]. Available: \url{https://mathworld.wolfram.com/LHospitalsRule.html}
\bibitem{mit_algorithms} E. Demaine, J. Ku, and J. Solomon, ``6.006 Introduction to Algorithms: Lecture Notes,'' MIT OpenCourseWare, 2020. [Online]. Available: \url{https://ocw.mit.edu/courses/6-006-introduction-to-algorithms-spring-2020/pages/lecture-notes/}
\end{thebibliography}

\end{document}
